% Source: https://artofproblemsolving.com/community/c3h2650413_sigma_mathematics_competition_smc_problems
% Stable identity: aops-topic-2650413-q10
Let $f(k)=\left\lfloor\sqrt k\right\rfloor$. For $j<k$ the slope is $\dfrac{(k-j)+f(k)-f(j)}{k-j}=1+\dfrac{f(k)-f(j)}{k-j}$, so $1<s<2$ exactly when $0<f(k)-f(j)<k-j$.

The values of $f$ form blocks: $f=1,2,3,4,5$ on $\{1,2,3\}$, $\{4,\dots,8\}$, $\{9,\dots,15\}$, $\{16,\dots,24\}$, $\{25\}$, of sizes $3,5,7,9,1$.
\begin{itemize}
\item $f(k)=f(j)$ (slope $1$) for $\dbinom32+\dbinom52+\dbinom72+\dbinom92=3+10+21+36=70$ pairs.
\item Since $f$ rises by at most $1$ per step, $f(k)-f(j)\le k-j$, with equality only if every step from $j$ to $k$ is a jump. As every block except the last has at least $3$ elements, this happens only for $(j,k)=(3,4),(8,9),(15,16),(24,25)$ (slope $2$).
\end{itemize}
So $\dbinom{25}{2}-70-4=226$ of the $300$ pairs work. The probability is $\dfrac{226}{300}=\dfrac{113}{150}$, and $m+n=\boxed{263}$.

\begin{asy}
usepackage("amsmath");
size(5cm);
defaultpen(fontsize(10pt));
for(int k=1;k<=25;++k){int f=floor(sqrt(k));dot((k,k+f),(f%2==0)?black:red);}
draw((0,0)--(27,0),gray,Arrow(4));draw((0,0)--(0,32),gray,Arrow(4));
draw((3,4)--(4,6),dashed);draw((8,10)--(9,12),dashed);draw((15,18)--(16,20),dashed);draw((24,28)--(25,30),dashed);
label("$y=x+1$",(3,4),NW,fontsize(8pt));label("$y=x+5$",(25,30),NW,fontsize(8pt));
\end{asy}

(Points with even and odd $\left\lfloor\sqrt k\right\rfloor$ are shown in black and red; the dashed segments are the four pairs of slope $2$.)
